\section{Hypothesis Testing}

In hypothesis testing, we assume an initial premise (usually concerning the value of an estimator) called the \emph{null hypothesis}, and we attempt to determine whether or not it is true by applying statistical techniques.

We also have another premise called the \emph{alternative hypothesis}, which is the negation of the null hypothesis.

\subsection{Null vs. Alternative Hypothesis}

When conducting quantitative \emph{research} to calibrate the value of an estimator, the \emph{known value} of the parameter is taken as the \emph{null hypothesis}, while the \emph{new value} found (from the research) is taken as the \emph{alternative hypothesis}.

In our case (finding the average age of our city), a researcher might claim that the age is \emph{less than 35}. This can serve as the \emph{null hypothesis}.

If a new agency claims that it is \emph{greater than 35}, then that can be designated as the \emph{alternative hypothesis}.

\subsection{Formal Notation}

We denote the hypotheses as:
\begin{itemize}
	\item $H_0$: Null hypothesis (the claim initially assumed to be true)
	\item $H_a$ or $H_1$: Alternative hypothesis (the claim that contradicts $H_0$)
\end{itemize}

\begin{ejemplo}
	A cereal manufacturer claims that the average weight of its boxes is 500 grams. A quality inspector suspects the weight is lower.
	
	\begin{align}
		H_0: \mu &= 500 \text{ g} \\
		H_a: \mu &< 500 \text{ g}
	\end{align}
\end{ejemplo}

\subsection{Types of Hypothesis Tests}

Depending on the alternative hypothesis, tests can be:

\begin{enumerate}
	\item \textbf{Left-tailed test:}
	\begin{align}
		H_0: \mu &= \mu_0 \\
		H_a: \mu &< \mu_0
	\end{align}
	
	\item \textbf{Right-tailed test:}
	\begin{align}
		H_0: \mu &= \mu_0 \\
		H_a: \mu &> \mu_0
	\end{align}
	
	\item \textbf{Two-tailed test:}
	\begin{align}
		H_0: \mu &= \mu_0 \\
		H_a: \mu &\neq \mu_0
	\end{align}
\end{enumerate}

\subsection{Errors in Hypothesis Testing}

When making decisions based on sample data, we can commit two types of errors:

\begin{definicion}[Type I Error]
	Rejecting the null hypothesis when it is actually true. The probability of committing this error is denoted by $\alpha$ (significance level):
	\begin{align}
		\alpha = P(\text{Reject } H_0 \mid H_0 \text{ is true})
	\end{align}
\end{definicion}

\begin{definicion}[Type II Error]
	Failing to reject the null hypothesis when it is actually false. The probability of committing this error is denoted by $\beta$:
	\begin{align}
		\beta = P(\text{Do not reject } H_0 \mid H_0 \text{ is false})
	\end{align}
\end{definicion}

\begin{observacion}
	The \emph{power} of a test is the probability of correctly rejecting a false null hypothesis:
	\begin{align}
		\text{Power} = 1 - \beta = P(\text{Reject } H_0 \mid H_0 \text{ is false})
	\end{align}
\end{observacion}

\subsection{Power and Sample Size}

The \emph{power} of a test is not a fixed value: it depends on the real effect size to be detected ($\mu_a - \mu_0$), on the sample size $n$, and on the chosen significance level $\alpha$. In experimental design, the researcher can fix both $\alpha$ and the desired power $1-\beta$ in advance to determine how many observations $n$ must be collected.

\begin{teorema}[Sample size for a target power, one-tailed $Z$-test]
	Consider a hypothesis test about the mean of a normal population with known standard deviation $\sigma$, of the form $H_0: \mu = \mu_0$ versus $H_a: \mu > \mu_0$ (or its left-tailed analogue). If the test is required to achieve a power $1-\beta$ to detect a real effect of size $\mu_a \neq \mu_0$ at significance level $\alpha$, the minimum required sample size is:
	\begin{align}
		n = \left( \frac{(Z_\alpha + Z_\beta)\, \sigma}{\mu_a - \mu_0} \right)^2,
		\label{eq:tam_muestra_potencia_en}
	\end{align}
	where $Z_\alpha$ and $Z_\beta$ are the upper standard normal quantiles corresponding to $\alpha$ and $\beta$ respectively. For a two-tailed test, $Z_\alpha$ is replaced by $Z_{\alpha/2}$.
\end{teorema}

\begin{observacion}
	Formula \eqref{eq:tam_muestra_potencia_en} reveals fundamental experimental-design relationships:
	\begin{itemize}
		\item The larger the desired effect size ($\mu_a - \mu_0$), the \textbf{smaller} the required sample size: large effects are easier to detect.
		\item The higher the required power (for example, moving from $1-\beta = 0.80$ to $1-\beta = 0.95$), the \textbf{larger} the required sample size, since $Z_\beta$ grows.
		\item The smaller $\alpha$ (a stricter test against false positives), the \textbf{larger} the required sample size as well, since $Z_\alpha$ grows.
	\end{itemize}
	This formalizes the intuition that there is a triangular trade-off among Type I error, Type II error, and the cost (sample size) of an experiment.
\end{observacion}

\subsection{Significance Level}

The \emph{significance level} $\alpha$ is the maximum acceptable probability of committing a Type I error. Common values are:
\begin{itemize}
	\item $\alpha = 0.10$ (10\%) - Low significance level
	\item $\alpha = 0.05$ (5\%) - Standard significance level
	\item $\alpha = 0.01$ (1\%) - High significance level
\end{itemize}

\begin{observacion}
	The choice of $\alpha$ depends on the context and the consequences of committing a Type I error. In medical or safety applications, stricter significance levels ($\alpha = 0.01$ or lower) are typically used.
\end{observacion}

\subsection{$p$-value}

The \emph{$p$-value} is the probability of obtaining a test result at least as extreme as the one observed, assuming that the null hypothesis is true.

\begin{definicion}[$p$-value]
	The $p$-value is:
	\begin{align}
		p\text{-value} = P(\text{Statistic} \geq \text{observed value} \mid H_0 \text{ is true})
	\end{align}
	for right-tailed tests, or the corresponding area depending on the type of test.
\end{definicion}

\subsection{Decision Rule}

The decision rule based on the $p$-value is:

\begin{itemize}
	\item If $p\text{-value} \leq \alpha$: \textbf{Reject} $H_0$
	\item If $p\text{-value} > \alpha$: \textbf{Do not reject} $H_0$
\end{itemize}

\begin{observacion}
	``Do not reject $H_0$'' is not the same as ``accept $H_0$''. It simply means there is not enough evidence to reject it.
\end{observacion}

\subsection{Complete Example}

A company claims that the average delivery time for its products is 3 days. A customer suspects the time is longer and takes a sample of $n = 36$ deliveries, obtaining $\bar{x} = 3.5$ days with $s = 1.2$ days.

\textbf{Step 1: Formulate hypotheses}
\begin{align}
	H_0: \mu &= 3 \text{ days} \\
	H_a: \mu &> 3 \text{ days}
\end{align}

\textbf{Step 2: Choose significance level}
\begin{align}
	\alpha = 0.05
\end{align}

\textbf{Step 3: Calculate test statistic}

Since $n = 36 \geq 30$, we use the $Z$-test:
\begin{align}
	Z = \frac{\bar{x} - \mu_0}{s/\sqrt{n}} = \frac{3.5 - 3}{1.2/\sqrt{36}} = \frac{0.5}{0.2} = 2.5
\end{align}

\textbf{Step 4: Calculate $p$-value}
\begin{align}
	p\text{-value} = P(Z > 2.5) = 1 - \Phi(2.5) \approx 0.0062
\end{align}

\textbf{Step 5: Make decision}

Since $p\text{-value} = 0.0062 < 0.05 = \alpha$, we reject $H_0$.

\textbf{Conclusion:} There is sufficient statistical evidence at the 5\% significance level to claim that the average delivery time is greater than 3 days.

\subsection{Interpretation of Results}

It is crucial to correctly interpret the results:

\begin{itemize}
	\item \textbf{Reject $H_0$:} There is sufficient statistical evidence to support the alternative hypothesis.
	
	\item \textbf{Do not reject $H_0$:} There is not sufficient statistical evidence to support the alternative hypothesis. This does not mean $H_0$ is true, only that we could not prove otherwise.
\end{itemize}

\subsection{Important Considerations}

\begin{enumerate}
	\item \textbf{Sample size:} Larger samples provide greater statistical power to detect real differences.
	
	\item \textbf{Statistical vs. practical significance:} A result can be statistically significant without being practically relevant.
	
	\item \textbf{Assumptions:} Hypothesis tests assume that certain underlying conditions (normality, independence, etc.) are met.
	
	\item \textbf{Multiple testing:} Performing many tests increases the probability of committing at least one Type I error.
\end{enumerate}

\subsection{Summary of the Procedure}

The general procedure for a hypothesis test is:

\begin{enumerate}
	\item Formulate $H_0$ and $H_a$
	\item Choose the significance level $\alpha$
	\item Select the appropriate test statistic
	\item Calculate the value of the statistic using the sample data
	\item Determine the $p$-value or critical region
	\item Make a decision: reject or do not reject $H_0$
	\item Interpret the result in the context of the problem
\end{enumerate}

In the following sections we will study specific tests for means ($Z$ and $t$ tests), proportions, and other common situations.
